\documentclass[10pt]{article}

\usepackage[utf8]{inputenc}

\usepackage[T1]{fontenc}

\usepackage{amsmath}

\usepackage{amsfonts}

\usepackage{amssymb}

\usepackage[version=4]{mhchem}

\usepackage{stmaryrd}

\usepackage{graphicx}

\usepackage[export]{adjustbox}

\graphicspath{ {./images/} }

\usepackage{bbold}



\begin{document}

в множестве $S_{j}$ задается расстояние



\[

d(f, g)=\sup _{x \in E^{n}}\left(|f x-g x|+\left\|d f_{x}-d g_{x}\right\|\right) .

\]



При всяком $j \in S$ в пространстве $S_{j} \times E^{n}$ имеется всюду плотное множество $D_{j}$ типа $G_{\delta}$, обладающее свойством: для всякого $(f, x) \in D_{j}$ точка $x$ условно экспоненциально устойчива относительно диффеоморфизма $f$ с индексом



\[

\operatorname{dim}\left\{\mathfrak{g} \in T_{x} E^{n}: \varlimsup_{m \rightarrow+\infty} \frac{1}{m} \ln \left|d f^{m} \mathfrak{g}\right|<0\right\},

\]



т. е. с индексом, равным числу отрицательных Ляпунова характеристических показателей уравнения в вариациях .



\begin{enumerate}

  \setcounter{enumi}{1}

  \item Для динамич. систем, заданных на замкнутом дифференцируемом многообразии, аналогичная теорема формулируется проще и является диффференциальнотопологическіІ инвариантной. Пусть Г'n-замкнутое дпффференцируемое многообразие. Множество $S$ всех дифффеоморфизмов $f$ класса $C^{1}$, отображающих $V^{n}$ на $V^{n}$, наделяется $C^{1}$-топологией. В пространстве $S \times V^{n}$ имеется всюду плотное множество $D$ типа $G_{\delta}$, обладаюшее свойством: для всякого $(f, x) \in D$ точка $x$ условно экспоненциально устойчива относительно диффеоморфизма $f$ с индексом

\end{enumerate}



\[

k(x)=\operatorname{dim}\left\{\mathfrak{x} \in T_{x} V^{n}: \varlimsup_{m \rightarrow+\infty} \frac{1}{m} \ln \left|d f^{m} \mathfrak{x}\right|<0\right\} .

\]



\begin{enumerate}

  \setcounter{enumi}{2}

  \item Для всякого диффеоморфизма $f: V^{n} \longrightarrow V^{n}$ замкнутого диффференци руемого многообразия $V^{n}$ для всякого инвариантного относительно $f$ распределения вероятностей на $V^{n}$ (б-алгебра к-рого содержит все борелевские множества) множество точек $x \in V^{n}$, условно экспоненциально устойчивых с индексом (4) относительно диффеоморфизма $f$, имеет вероятность 1.

\end{enumerate}



\begin{center}

\includegraphics[max width=\textwidth]{2023_07_09_fe214dbe9e525639254fg-1(1)}

\end{center}



\includegraphics[max width=\textwidth, center]{2023_07_09_fe214dbe9e525639254fg-1}

нова и ее приложения к вопросам' устойчивости, М., 1966; [3] Д е м и д о в и ч Б. П., Лекции по математической теории устойчивости, М., 1967, [4] И з о б о в Н. А., в кн.: Итоги науки и техники. Математический анализ, т. 12, М., 1974 , с. $71-146$; с. $55-112$.

В. М. Миллионциков.



УСЛОВНО ПЕРИОДИЧЕСКАЯ ФУНКЦИЯ - Функция $A \circ \varphi$, являющаяся композицией $2 \pi-$ приодич. функции $A: T^{n} \longrightarrow \mathbb{C}$, где $T^{n}$ есть $n$-мерный тор, и функции $\varphi: \mathbb{R} \longrightarrow \mathbb{R} n$ такой, что $\dot{\varphi}=\omega$, где $\omega=\left(\omega_{1}, \ldots, \omega_{n}\right)-$ постоянный вектор $c$ рационально линейно независимыми компонентами. Примером У. П. ф. служит отрезок ряда Фурье



\[

\sum_{i=1}^{n}\left[A_{i} \sin \left(\omega_{i} t+\psi_{i}\right)+B_{i} \cos \left(\omega_{i} t+\psi_{i}\right)\right]

\]



где



\[

A:=\sum_{i=1}^{n}\left[A_{i} \sin \varphi_{i}+B_{i} \cos \varphi_{i}\right]

\][



\textbackslash varphi:=\textbackslash left(\textbackslash varphi\_\{1\}(t), ···, \textbackslash varphi\_\{n\}(t)\textbackslash right)=\textbackslash left(\textbackslash omega\_\{1\} t+\textbackslash psi\_\{1\}, ···, \textbackslash omega\_\{n\} t+\textbackslash psi\_\{n\}\textbackslash right) .

]



ЕсліІ У. П. Ф.- непрерывная функция, то она совпадает с көазипериодической функцией с периодами $\omega_{1}, \ldots, \omega_{n}$. П. В. Комленко.



УСЈОВНО ПЕРИОДИЧЕСКОЕ ДВИЖЕНИЕ - прямолинейное движение материальної точки, закон к-рого выражается деӥствительной условно периодической фуннuиeй.



УСЛОВНО ПОЛНАЯ РЕIIIEТКА - репетка, в к-рої каждое непустое ограниченное подмножество имеет точную верхнюю грань и точную нижнюю грань. Примером У. ю. р. может служить множество всех действительных чисел с обычным порядком. T. C. Фоданова.



УСЛОВНОЕ МАТЕМАТИЧЕСКОЕ ОЖИДАНИЕ с Л Уч а й н о й в е л и ч и н -фурукция элементарного собы- тия, характеризующая случайную величину по отнопению к нек-рой $\sigma$-алгебре. Пусть $(\Omega, \mathcal{A}, \mathrm{P})-$ вероятностное пространство, $X$ - заданная на нем слутайная величина с конечным математич. ожиданием, $\mathfrak{B}$ есть $\sigma$-алгебра, $\mathfrak{B} \subseteq \mathcal{A}$. У. м. о. случайной велитины $X$ относительно $\sigma$-алгебры $\mathfrak{B}$ наз. случайная величина $E(X \mid \mathfrak{B})$, измеримая относительно б-алгебры $\mathfrak{B}$ и такая, что



\[

\int_{B} X P(d \omega)=\int_{B} E(X \mid B) P(d \omega)

\]



для каждого $B \in \Re$. Если математич. ожидание случайной величины $X$ бесконечно (но определено), т. е. конечна только одна из величин $\mathrm{E} X^{+}=\mathrm{E} \max (0, X)$ и $\mathrm{EX}-=-\operatorname{Emin}(0, X)$, то определение У. м. о. посредством (*) имеет смысл, но $\mathrm{E}(X \mid \mathfrak{B})$ может принимать бесконетные значения.



У. м. о. определяется однозначно с точностью до эквивалентности. В отличие от математического ожидания, являющегося числом, У. м. о. представляет собой фуункцию (случайную величину).



Свойства У. м. о. аналогичны свойствам математич. ожидания:



\begin{enumerate}

  \item $\mathrm{E}\left(X_{1} \mid \mathfrak{B}\right) \leqslant \mathrm{E}\left(X_{2} \mid \mathfrak{B}\right)$, если почти наверное $X_{1}(\omega) \leqslant X_{2}(\omega)$



  \item $\mathrm{E}(c \mid \mathfrak{B})=c$ для любого действительного $c$;



  \item $\mathrm{E}\left(\alpha X_{1}+\beta X_{2} \mid \mathfrak{B}\right)=\alpha \mathrm{E}\left(X_{1} \mid \mathfrak{B}\right)+\beta \mathrm{E}\left(X_{2} \mid \mathfrak{B}\right)$ для любых действительных $\alpha$ и $\beta$;



  \item $|\mathrm{E}(X \mid \mathfrak{B})| \leqslant \mathrm{E}(|X| \mid \mathfrak{B})$;



  \item $g(E(X \mid \mathfrak{B})) \leqslant E(g(X) \mid \mathfrak{B})$ для выпуклых функций $g(x)$.



\end{enumerate}



Кроме того, имеют место следующие специфические для У. м. о. свойства:



\includegraphics[max width=\textwidth, center]{2023_07_09_fe214dbe9e525639254fg-1(2)}

$\mathrm{E}(X \mid \mathfrak{B})=\mathrm{EX}$



\begin{enumerate}

  \setcounter{enumi}{6}

  \item $\mathrm{E}(X \mid \mathcal{A})=X$



  \item $\mathrm{E}(\mathrm{E}(X \mid \mathfrak{B}))=\mathrm{E} X$



  \item если $X$ не зависит от $\sigma$-алгебры $\mathfrak{B}$, то $E(X \mid \mathfrak{B})=$ $=\mathrm{E} \cdot \boldsymbol{X}$



  \item если $Y$ измерима относительно б-алгебры $\mathfrak{B}$, то $\mathrm{E}(X Y \mid \mathfrak{B})=Y \mathrm{E}(X \mid \mathfrak{B})$.



\end{enumerate}



Имеет место теорема о сходимости под знаком У . м. о.: если $X_{1}, X_{2}, \ldots$-последовательность случайных велитин, $\left|X_{n}\right| \leqslant Y, n=1,2, \ldots, \mathrm{E} Y<\infty$ и $X_{n} \longrightarrow X$ почти наверное, то почти наверное $\mathrm{E}\left(X_{n} \mid \mathfrak{B}\right) \longrightarrow \mathrm{E}(X \mid \mathfrak{B})$.



У. м. о. случайной величины $X$ относительно случайной величины $Y$ определяется как У. м. о. $X$ относительно $\sigma$-алгебры, порожденной $Y$.



Частным случаем У. м. о. является условная вероятHocmb.



Лит.: [1] К о лм огоро в А. Н., Основные понятия теарии вероятностей, 2 изд., М., $1974 ;$ [2] П р о х о р о в Ю. В., [3] Н е в ё Ж., Математические основы теории вероятностей, пер. с франц., М., 1969, [4] JЈ о ә в М., Теория вероятностей, пер. с англ., М., 1962.



УСЛОВНЗОЕ РАСПРЕДЕЛЕНИЕ - Функция олементарного события и борелевского множества, при каждом фикси рованном әлементарном событии являющаяся päcпределением вероятностей, а при каждом фиксированном борелевском множестве - условной вероятностью.



Пусть $(\Omega, \mathcal{A}, \mathrm{P})$ - вероятностное пространство, $\mathfrak{B}$ есть $\sigma$-алгебра борелевских множеств на прямой, $X-c л у-$ чайная неличина, определенная на $(\Omega, \mathcal{A}), \mathfrak{F}$-под$\sigma$-алгебра $\mathcal{A}$. Функция $Q(\omega, B)$, определенная на $\Omega \times \mathfrak{B}$, наз. (регу ля р ны м) ус ловны м р ас п р едел ени е случайной величины $X$ относительно б-алгебры $\mathfrak{F}$, если:



a) при фиксированном $B \in \mathfrak{B}$ функция $Q(\omega, B)$ F-измерима,



б) с вероятностью единица при фиксированном $\omega$ функция $Q(\omega, B)$ является вероятностной мерой на $\mathfrak{B}$,





\end{document}